\clearpage\thispagestyle{empty}
\pdfbookmark[0]{Quellen zu Beobachtung der Mondfinsternis vom 31. Mai 1844 in Göttingen}{lunareclipse1844-sources}
\begin{center}\textbf{Quellen zu Beobachtung der Mondfinsternis vom 31. Mai 1844 in Göttingen}\end{center}
\noindent Diese KI-Übertragung ins heutige Deutsch folgt dem Bericht auf S.~466 der \href{https://gdz.sub.uni-goettingen.de/id/PPN236007467}{\textit{Werke}, Band~VI} (Göttingen, 1874). Carl Friedrich Gauß und Goldschmidt machten die Beobachtungen. Gauß unterzeichnete den Bericht. Er erschien in den \textit{Astronomischen Nachrichten}, Band~XXII, Nr.~506, \href{https://api.digitale-sammlungen.de/iiif/image/v2/bsb10049447_00024/full/1800,/0/default.jpg}{Sp.~31--32}, ausgegeben am 20. Juni 1844. Der Katalog der Bayerischen Staatsbibliothek datiert den gebundenen Band auf 1845.
\par\smallskip\noindent Beim Eintritt von Aristarch gibt die Werkausgabe für Goldschmidts Beobachtung 32 Minuten an, der Erstdruck 52. Im Erstdruck ist zudem vom Rand des Mondes die Rede, der am weitesten vom Mittelpunkt des Schattens entfernt ist. Die Werkausgabe spricht vom entfernten Rand. Die Übertragung folgt der Werkausgabe und belässt Goldschmidts 15~1~21 in der Zeile Plinius~I.
\par\smallskip\noindent Der kurze Strich vor Goldschmidts Eintrag zum Eintritt von Dionysius ist als Minuszeichen gesetzt. Gauß’ erste Austrittszeit behält den Punkt nach dem Stundenzeichen. Die Punktreihe beim Austritt von Grimaldi~II behält ihr Komma. KI-Übertragung und KI-Satz.
\clearpage
